\begingroup
\renewcommand{\arraystretch}{1.45}
\Needspace{22\baselineskip}
\InterpolationArticle{26}
Wäre die Methode der Artikel 20, 22 auf alle $\pi$ vorgegebenen Werte unmittelbar angewandt worden, ohne sie zuvor in $\nu$ kleinere Perioden zu zerlegen, so wäre eine Funktion ($X'''$) der $\tfrac12\pi-\tfrac12^{\mathrm{ten}}$ oder $\tfrac12\pi^{\mathrm{ten}}$ Ordnung entstanden, je nachdem, ob $\pi$ ungerade oder gerade ist; diese werden wir nun mit der Funktion $X''$ vergleichen. Da die letztere, wenn sie gehörig reduziert ist, offenbar bis zur $\mu n+m^{\mathrm{te}}$ Ordnung reicht, werden wir drei Fälle unterscheiden.

I. Sind sowohl $\mu$ als auch $\nu$ ungerade, so wird $\mu n+m=\tfrac12\mu\nu-\tfrac12=\tfrac12\pi-\tfrac12$ sein. Hieraus ist klar, dass $X''$ mit $X'''$ vollständig identisch sein muss.

\Needspace{31\baselineskip}
II. Ist $\mu$ gerade und $\nu$ ungerade, so ergibt sich $\mu n+m=\tfrac12\mu\nu=\tfrac12\pi$, und somit ist $X''$ von dersel\SourceLeaf{314}%
ben Ordnung wie $X'''$. Daher genügt es, die letzten Glieder beider Funktionen zu vergleichen. Es ist nicht schwer zu sehen, dass diese letzten Glieder in $X''$ ausschließlich aus den Gliedern $\gamma^m\cos mx+\delta^m\sin mx$ entstehen, und zwar aus denjenigen Gliedern der Koeffizienten $\gamma^m,\delta^m$, welche $\cos ny$ und $\sin ny$ oder $\cos\mu nx$ und $\sin\mu nx$ enthalten. Diese Glieder seien
\[
\begin{aligned}
&\text{in }\gamma^m\quad.\,.\,.\,.\quad k\cos\mu nx+l\sin\mu nx\\[5pt]
&\text{in }\delta^m\quad.\,.\,.\,.\quad k'\cos\mu nx+l'\sin\mu nx
\end{aligned}
\]
Hieraus entsteht in $X''$
\[
\begin{aligned}
&(k\cos\mu nx+l\sin\mu nx)\cos mx\\[5pt]
&\qquad{}+(k'\cos\mu nx+l'\sin\mu nx)\sin mx
\end{aligned}
\]
woraus die folgenden Glieder der $\mu n+m^{\mathrm{ten}}$ Ordnung hervorgehen:
\[
\begin{aligned}
&\tfrac12(k-l')\cos(\mu n+m)x+\tfrac12(l+k')\sin(\mu n+m)x
\end{aligned}
\]
Nun lässt sich zwar ohne Schwierigkeit \textit{aus der Bildung der Koeffizienten} zeigen, dass
\[
\begin{aligned}
&(k-l')\sin(\mu n+m)a=(l+k')\cos(\mu n+m)a
\end{aligned}
\]
\Needspace{46\baselineskip}
doch ziehen wir die folgende Methode vor, die eleganter erscheint. Aus Artikel 24 folgt, dass
\[
\begin{aligned}
&\gamma^m\sin\tfrac12y-\delta^m\cos\tfrac12y=0
\end{aligned}
\]
\textit{sowohl} dann gilt, wenn für $y$ der Wert $\mu a$ und für $\gamma^m,\delta^m$ die aus der ersten Periode hergeleiteten Werte eingesetzt werden,
\textit{als auch} dann, wenn für $y$ der Wert $\mu a'$ und für $\gamma^m,\delta^m$ die aus der zweiten Periode hergeleiteten Werte eingesetzt werden, usw. Werden also für $\gamma^m,\delta^m$ die gefundenen unbestimmten Ausdrücke eingesetzt, sodass $\gamma^m\sin\tfrac12y-\delta^m\cos\tfrac12y$ zu einer Funktion des Winkels $y$ von der Form $G$ wird, die bis zur $n+\tfrac12=\tfrac12\nu$ Ordnung reicht, so wird diese für alle $\nu$ Werte von $y$, nämlich $\mu a,\mu a',\mu a''$ usw., $=0$ und ist daher nach Artikel 14 durch
\[
\begin{aligned}
&\sin\tfrac12(y-\mu a)\sin\tfrac12(y-\mu a')\sin\tfrac12(y-\mu a'')\text{ usw.}
\end{aligned}
\]
teilbar, somit auch durch $\sin\tfrac12\nu(y-\mu a)$, und folglich von der Form $h\sin\tfrac12\nu(y-\mu a)$, wobei $h$ eine feste Größe ist. Daher müssen die letzten Glieder der entwickelten Funktion $\gamma^m\sin\tfrac12y-\delta^m\cos\tfrac12y$
\[
\begin{aligned}
&-h\sin\tfrac12\pi a\cos\tfrac12\nu y+h\cos\tfrac12\pi a\sin\tfrac12\nu y
\end{aligned}
\]
sein und alle vorangehenden verschwinden; jene letzten Glieder werden aber offenbar
\SourceLeaf{315}
\[
\begin{aligned}
&=-\tfrac12(l+k')\cos(n+\tfrac12)y+\tfrac12(k-l')\sin(n+\tfrac12)y
\end{aligned}
\]
Somit wird
\[
\begin{aligned}
&\tfrac12(k-l')=h\cos\tfrac12\pi a=h\cos\tfrac12(\mu n+m)a\\[5pt]
&\tfrac12(l+k')=h\sin\tfrac12\pi a=h\sin\tfrac12(\mu n+m)a\qquad\text{Q.~E.~D.}
\end{aligned}
\]
sein. Diese Gleichung wird aber genau das Kriterium für die letzten Koeffizienten der Funktion
$X'''$ sein, wie in Artikel 24 gezeigt worden ist; daher wird auch in diesem Fall die Funktion $X''$ mit $X'''$
vollständig identisch sein.

\Needspace{32\baselineskip}
III. Ist $\nu$ gerade, so ergibt sich $\mu n+m=\tfrac12\pi+m$, und somit ist $X''$ von höherer Ordnung
als $X'''$. Es ist aber leicht zu sehen, dass alle Glieder in $X''$, welche die
$\tfrac12\pi^{\mathrm{te}}$ Ordnung überschreiten, ausschließlich aus denjenigen Gliedern der einzelnen
Koeffizienten $\gamma',\delta',\gamma'',\delta''$ usw. hervorgehen, welche $\cos ny$ und
$\sin ny$ oder $\cos\mu nx$ und $\sin\mu nx$ enthalten. Diese Glieder seien im Ausdruck für
einen beliebigen Koeffizienten $\gamma^\lambda$ oder $\delta^\lambda$ gleich $k\cos\mu nx+l\sin\mu nx$;
aus ihnen entsteht in $X''$
\[
\begin{aligned}
&(k\cos\mu nx+l\sin\mu nx)\cos\lambda x\quad\text{im ersten Fall}\\[5pt]
&\text{oder}\quad(k\cos\mu nx+l\sin\mu nx)\sin\lambda x\quad\text{im zweiten Fall}
\end{aligned}
\]
und hieraus
\[
\begin{aligned}
&\tfrac12k\cos(\mu n+\lambda)x+\tfrac12l\sin(\mu n+\lambda)x\\[5pt]
&+\tfrac12k\cos(\mu n-\lambda)x+\tfrac12l\sin(\mu n-\lambda)x
\end{aligned}
\]
im ersten Fall und
\[
\begin{aligned}
&-\tfrac12l\cos(\mu n+\lambda)x+\tfrac12k\sin(\mu n+\lambda)x\\[5pt]
&+\tfrac12l\cos(\mu n-\lambda)x-\tfrac12k\sin(\mu n-\lambda)x
\end{aligned}
\]
im zweiten Fall. Nach Artikel 24 muss man aber für $\cos(\mu n+\lambda)x$
oder $\cos(\tfrac12\pi+\lambda)x$ den Ausdruck $\cos\pi a\cos(\tfrac12\pi-\lambda)x+\sin\pi a\sin(\tfrac12\pi-\lambda)x$ einsetzen,
und für $\sin(\mu n+\lambda)x$ oder $\sin(\tfrac12\pi+\lambda)x$ den Ausdruck
$\sin\pi a\cos(\tfrac12\pi-\lambda)x-\cos\pi a\sin(\tfrac12\pi-\lambda)x$.
Daher werden wir im ersten Fall
\[
\begin{aligned}
&\tfrac12(k+k\cos\pi a+l\sin\pi a)\cos(\tfrac12\pi-\lambda)x\\[5pt]
&+\tfrac12(l+k\sin\pi a-l\cos\pi a)\sin(\tfrac12\pi-\lambda)x
\end{aligned}
\]
und im zweiten Fall
\[
\begin{aligned}
&\tfrac12(l-l\cos\pi a+k\sin\pi a)\cos(\tfrac12\pi-\lambda)x\\[5pt]
&-\tfrac12(k+l\sin\pi a+k\cos\pi a)\sin(\tfrac12\pi-\lambda)x
\end{aligned}
\]
\SourceLeaf{316}
erhalten. Aus der Gleichung $k\sin\tfrac12\pi a=l\cos\tfrac12\pi a$ aber, die in beiden Fällen
zwischen $k$ und $l$ gelten muss (Artikel 24), folgt ohne Schwierigkeit
\[
\begin{aligned}
&\tfrac12(k+k\cos\pi a+l\sin\pi a)=k\\[5pt]
&\tfrac12(l+k\sin\pi a-l\cos\pi a)=l
\end{aligned}
\]
\Needspace{32\baselineskip}
Hieraus folgt schließlich, dass man für diejenigen Teile in $X''$, die aus den letzten Gliedern
\[
\begin{aligned}
&k\cos ny+l\sin ny
\end{aligned}
\]
eines jeden Koeffizienten $\gamma^\lambda$ oder $\delta^\lambda$ hervorgehen würden,
wobei $\lambda$ einen Index größer als 0 bezeichnet, unmittelbar
\[
\begin{aligned}
&k\cos(\tfrac12\pi-\lambda)x+l\sin(\tfrac12\pi-\lambda)x\\[5pt]
&\qquad\text{im ersten Fall, oder}\\[5pt]
&l\cos(\tfrac12\pi-\lambda)x-k\sin(\tfrac12\pi-\lambda)x\\[5pt]
&\qquad\text{im zweiten Fall}
\end{aligned}
\]
einsetzen darf. Auf diese Weise erhalten wir anstelle der Funktion $X''$, die bis zur $\tfrac12\pi+m^{\mathrm{te}}$
Ordnung reichen würde, eine andere, $X''''$, welche die $\tfrac12\pi^{\mathrm{te}}$
Ordnung nicht überschreitet. Es bleibt daher nur noch, die Koeffizienten des Kosinus und Sinus des Winkels
$\tfrac12\pi x$, die in dieser Funktion offenbar dieselben geblieben sind wie in $X''$, mit
den entsprechenden Koeffizienten in $X'''$ zu vergleichen. Bei einiger Aufmerksamkeit
sieht man aber, dass diese Glieder ausschließlich aus dem ersten Koeffizienten $\gamma$ entstehen, und zwar
aus dessen letzten Gliedern. Werden diese durch $k\cos ny+l\sin ny$ dargestellt, so werden jene
$k\cos\tfrac12\pi x+l\sin\tfrac12\pi x$ und sind daher notwendig mit den entsprechenden Gliedern in
$X'''$ identisch. Hieraus folgt also, dass die Funktion $X''''$, die nach den
soeben gegebenen Vorschriften anstelle der Funktion $X''$ erhalten wird, mit der Funktion $X'''$
vollständig identisch wird.

\endgroup
